\chapter{Introduction}
\label{chap:intro}

%% ──────────────────────────────────────────────────────────────────────────
\section{Motivation}
\label{sec:intro_motivation}

The ability to generate or complete three-dimensional scenes is one of the
more ambitious goals in computer vision and graphics. Applications span
rapid content creation, embodied-AI simulation, autonomous navigation testing,
and robot planning in unseen environments. Progress towards this goal has
accelerated in recent years, partly because the representation of 3D scenes has
matured considerably.

3D Gaussian Splatting~\citep{kerbl2023gaussian} has become the dominant
representation for photorealistic scene reconstruction. A scene is modelled as
a set of anisotropic Gaussian splats, each carrying a position, an opacity, a
colour, and an anisotropic covariance parameterised by a rotation and a scale.
The representation renders in real time, matches or exceeds the quality of
earlier implicit approaches~\citep{mildenhall2020nerf}, and is explicitly
structured, meaning a network can in principle produce the primitives
directly rather than inferring them from a coordinate query.

Yet despite this progress in representation, scene-level 3DGS \emph{generation}
remains largely unexplored. The standard recipe for generation in neighbouring
domains is well-established: first compress the raw signal into a smooth,
compact latent space with a variational autoencoder~\citep{kingma2014vae},
then train a generative model on that latent~\citep{rombach2022ldm,
ho2020ddpm,lipman2023flowmatching}. This recipe has been proven at the object
level for 3D content~\citep{zeng2022lion,zhang2024gaussiancube,zhao2023michelangelo},
but scaling it to whole scenes is a different problem. A scene contains
hundreds of thousands of unordered Gaussian primitives, spans an unbounded
volume, mixes dozens of object categories, and comes with no canonical
coordinate frame or ordering. Building the scene-level VAE that would underpin
this recipe is the central challenge this thesis addresses.

Can3Tok~\citep{gao2025can3tok} established for the first time that a
scene-level 3DGS VAE is feasible: it demonstrated convergence on small training
sets and showed that the resulting latent could be used for generation.
This was an important existence proof, but it left open a question that became
the organising challenge of this thesis: can such an autoencoder be scaled
to thousands of heterogeneous scenes, and what is required to achieve that?

During the early stages of this project, the answer to that question proved
to be \emph{no, not without significant changes}. Training the Can3Tok
architecture on thousands of pre-fitted, frozen 3DGS scenes failed to
converge to anything useful. Understanding why, and fixing it, is what this
thesis is about.

%% ──────────────────────────────────────────────────────────────────────────
\section{Problem Statement}
\label{sec:intro_problem}

This thesis studies the construction and training of a semantically aware,
scalable, scene-level 3DGS variational autoencoder as Stage~1 of a two-stage
generative pipeline for scene generation and completion. The encoder must map
a variable-size, unordered Gaussian scene into a fixed-size structured latent
$Z \in \bR^{T \times d}$, and the decoder must reconstruct the scene from
it. Downstream, a flow-matching Stage~2 model should be trainable on the
resulting latent distribution.

The central difficulty we identify is that the per-primitive parameters of a
\emph{pre-fitted} 3DGS scene are structurally ambiguous in two compounding
ways.

\paragraph{The ordering problem.} A 3DGS scene is a set; any permutation of
its elements represents the same scene. A decoder that is asked to reconstruct
``the $i$-th Gaussian'' is being asked to reconstruct a target that has no
stable identity from one scene to the next. Without a canonical ordering,
the slot-aligned reconstruction loss has no persistent signal and the model
cannot converge.

\paragraph{The placement and orientation problem.} Even within a small spatial
region, the exact position of any individual Gaussian is largely a degree of
freedom of the fitting optimiser. And the covariance
$\vsigma_i = R(q_i)\,\mathrm{diag}(s_i^2)\,R(q_i)^\top$ is invariant to
sign flips and permutations of the quaternion and scale, meaning a naive loss
on the raw quaternion $q_i$ penalises differences that do not change the actual
Gaussian shape.

Both problems become much harder at scale. At 300 scenes the model has enough
repetition to memorise individual scenes despite these ambiguities. At 3,000
or 6,000 scenes, memorisation is no longer possible and the problems become
the dominant obstacle.

There is also a subtler interaction that we discovered empirically: the
\emph{sampling budget}. In Can3Tok, the Gaussian splats are jointly optimised
alongside the autoencoder, so the fitting process can co-adapt to the
tokenisation. With pre-fitted, frozen splats there is no such co-adaptation.
As a result, the number of Gaussians per latent token (the ``block size''
$g = \lceil N/T \rceil$) directly controls how much structural ambiguity the
decoder must resolve per token, and using a smaller budget is substantially
better in this regime.

%% ──────────────────────────────────────────────────────────────────────────
\section{Research Questions}
\label{sec:intro_rqs}

The work addresses one main question and three sub-questions.

\medskip
\noindent\textbf{Main question}: \textit{What is required to scale a
scene-level 3DGS variational autoencoder from a few hundred training scenes
to several thousand heterogeneous scenes, and which design choices matter most?}

\medskip
\noindent\textbf{RQ1 (Ordering and identifiability).} Does imposing a canonical
spatial order on the Gaussian set restore convergence at scale, and how do
different ordering strategies compare?

\noindent\textbf{RQ2 (Architecture and generalisation).} Does replacing the
global encoder bottleneck with a spatially local encoder close the
train-to-test generalisation gap that appears at scale?

\noindent\textbf{RQ3 (Semantic structure).} Can per-primitive semantic
supervision and latent disentanglement be incorporated into the small-scale
regime without degrading reconstruction, motivating their inclusion in the
final design?

%% ──────────────────────────────────────────────────────────────────────────
\section{Contributions}
\label{sec:intro_contributions}

The contributions of this thesis are as follows.

\begin{enumerate}
\item \textbf{Identification of the scaling barriers.} We show, through
  systematic experimentation, that the ordering of the Gaussian set is the
  primary obstacle to scaling. Without a stable slot-to-position correspondence,
  the reconstruction target is not learnable at thousands of scenes.

\item \textbf{Fixed-frame Hilbert serialisation.} We apply a space-filling
  Hilbert curve in a fixed canonical coordinate frame to impose a consistent,
  locality-preserving order on the Gaussian set~\citep{skilling2004hilbert,
  wu2024ptv3}. We show that this is more effective than opacity ranking
  (the Can3Tok default) and Z-order (Morton) sorting.

\item \textbf{Anchor-relative position decoding.} We introduce per-token
  spatial anchors (block centroids) and decode each Gaussian's position as
  $\hat{\vmu}_i = a_{\tau(i)} + \rho\,\tanh(o_i)$, so the decoder predicts
  a bounded local offset rather than an absolute coordinate. This matches the
  Hilbert block structure and substantially reduces the within-block structural
  ambiguity the decoder must resolve.

\item \textbf{Improved reconstruction objectives.} We combine a Sinkhorn
  entropic-OT set loss~\citep{cuturi2013sinkhorn} that matches primitives
  within blocks by optimised assignment, with a Bures co-diagonalisation
  covariance loss that measures orientation and shape in a manner invariant to
  quaternion parameterisation.

\item \textbf{Spatially local encoder and token-local decoder.} We replace the
  globally-mixed bottleneck with a windowed cross-attention encoder and a
  shared per-token MLP decoder. We show that this is the single change most
  responsible for closing the train-to-test generalisation gap.

\item \textbf{Small-scale semantic exploration.} We conduct 300-scene
  experiments exploring per-primitive contrastive objectives and latent
  disentanglement, demonstrating that these are feasible and produce
  meaningful semantic structure in the latent.

\item \textbf{Stage-2 generative design.} We specify and partially implement
  a hierarchical flow-matching generator~\citep{lipman2023flowmatching,
  peebles2023dit} built on the Stage-1 latent, targeting both unconditional
  scene generation and partial-scan completion.
\end{enumerate}

We are upfront about scope: the large-scale ablation training is underway
at the time of writing and full quantitative results will be incorporated as
they become available. The Stage-2 training has not yet completed, and scene
completion remains the primary direction for future work.

%% ──────────────────────────────────────────────────────────────────────────
\section{Thesis Structure}
\label{sec:intro_structure}

Chapter~\ref{chap:related} positions this work in the literature.
Chapter~\ref{chap:method} presents the full method in detail.
Chapter~\ref{chap:results} presents the experimental findings: the 300-scene
exploratory results first, then the large-scale comparisons with
placeholders for in-progress runs. Chapter~\ref{chap:discussion} interprets
the findings and discusses limitations. Chapter~\ref{chap:conclusion}
answers the research questions. The appendix records all configuration details
for reproducibility.